\section{Secondary consequences of connected factorization}
\label{sec:supp-connected-consequences}

\begin{corollary}[Connected inversion]\label{cor:connected-inversion}
The connected birth counts are determined from the total counts by
\begin{equation}\label{eq:connected-inversion}
A^{\mathrm{conn}}(n)=
\sum_{\Pi\in\mathsf{Part}(\mathcal P(n))}
(-1)^{|\Pi|-1}(|\Pi|-1)!
\prod_{C\in\Pi}A(n_C).
\end{equation}
The same inversion holds for minimal generators:
\[
M^{\mathrm{conn}}(n)=
\sum_{\Pi\in\mathsf{Part}(\mathcal P(n))}
(-1)^{|\Pi|-1}(|\Pi|-1)!
\prod_{C\in\Pi}\#\mathcal M_{n_C}.
\]
\end{corollary}

\begin{proof}
Equation~\eqref{article-eq:partition-formula} is the ordinary exponential formula on
the finite set $\mathcal P(n)$ with block weight
$C\mapsto A^{\mathrm{conn}}(n_C)$. M\"obius inversion in the partition lattice
\cite[Section~3.7]{Stanley2012EnumerativeCombinatoricsI}
gives
\[
\mu_{\mathsf{Part}}(\Pi,\hat1)=(-1)^{|\Pi|-1}(|\Pi|-1)!,
\]
where $\hat1$ is the partition containing all coordinates in one block.
Applying the same partition-lattice inversion to
\eqref{article-eq:minimal-partition-formula} gives the minimal-generator formula.
\end{proof}

\begin{corollary}\label{cor:four-support-reduction}
If $\omega(n)=4$, every element of $B_n$ and every element of
$\mathcal M_n$ is obtained uniquely by choosing connected blocks whose sizes
form one of the set-partition types
\[
4,\qquad 3+1,\qquad 2+2,\qquad 2+1+1,\qquad 1+1+1+1.
\]
Consequently, after the one-, two-, and three-coordinate classifications are
known, the only data not reduced by this corollary are the connected
four-coordinate sets $B_n^{\mathrm{conn}}$ and
$\mathcal M_n^{\mathrm{conn}}$.
\end{corollary}

\begin{proof}
This is Theorems~\ref{article-thm:birth-connected-factorization} and
\ref{article-thm:mingenerator-connected-factorization} specialized to the set
$\mathcal P(n)$ of size four. Every disconnected partition has all blocks of
size at most three, so its factors are controlled by the previously proved
one-, two-, and three-coordinate classifications. The only partition with a
block of size four is the single-block partition, which is precisely the
connected four-coordinate case.
\end{proof}

\begin{example}[Symbolic support patterns through the four-coordinate reduction]
\label{ex:symbolic-support-patterns}
Identify the prime-coordinate set with $[k]$ and write
$\mathcal A_{I,J}$ for $\mathcal A_{I,J}(n)$. The following table gives a
compact reading of the witness-cover rules in small support. Each displayed
product is to be retained exactly when the indicated slots are nonempty and the
chosen atom representatives have distinct underlying primes.
\[
\begin{array}{ccl}
\toprule
\omega(n) & \text{symbolic cover} & \text{reading}\\
\midrule
1 & \theta\in\mathcal A_{\{1\},\{1\}} &
\text{one atom supplies the only coordinate}\\
2 & \theta\in\mathcal A_{\{1,2\},\{1,2\}}\ \text{or}\
\theta_1\theta_2,
\ \theta_i\in\mathcal A_{\{i\},\{i\}} &
\text{one full atom or two singleton blocks}\\
3 & \theta_{12}\theta_{23},\
\theta_{12}\in\mathcal A_{\{1,2\},\{1,2\}},\
\theta_{23}\in\mathcal A_{\{2,3\},\{2,3\}} &
\text{a connected three-coordinate row, e.g. }\Gamma_4\\
4 & m_Cm_{C'},\ C\sqcup C'=[4],\ |C|,|C'|\le3 &
\text{a nonconnected }3+1\text{ or }2+2\text{ assembly}\\
4 & m_Cm_i m_j\ \text{or}\ m_1m_2m_3m_4 &
\text{the }2+1+1\text{ and }1+1+1+1\text{ assemblies}\\
\bottomrule
\end{array}
\]
In the fourth and fifth rows, each factor $m_C$ is a connected minimal
generator for the coordinate block $C$, viewed inside $n_C$, and then embedded
back into $n$ by Theorem~\ref{article-thm:mingenerator-connected-factorization}. Thus
all nonconnected four-coordinate patterns are assembled from the already
classified one-, two-, and three-coordinate pieces.  The connected single-block
case $C=[4]$ is not classified here.
\end{example}

